\documentclass[10pt]{article}

\usepackage[utf8]{inputenc}

\usepackage[T1]{fontenc}

\usepackage{amsmath}

\usepackage{amsfonts}

\usepackage{amssymb}

\usepackage[version=4]{mhchem}

\usepackage{stmaryrd}



\begin{document}

связанная с Т. структура - структура близости. Понятие близости пространства. основано на отношении близости между подмножествами пространства- в отличие от понятия топологич. пространства.



В зависимости от широты класса исследуемых топологич. пространств меняется характер постановки основной задачи Т. Так, ограничиваясь узким классом пространств, ставят задачу различить их между собой в терминах топологич. инвариантов с точностью до гомеоморфизмов. Задача эта выглядит вполне естественной, напр., в рамках класса топологич. многообразий - но даже и здесь она оказывается чрезвычайно трудной и заведомо алгоритмически неразрешимой. Сложность задачи различения многообразий с точностью до гомеоморфизма приводит к необходимости рассмотреть более широкое, чем гомеоморфность, отношение г о мо топич. ә к в и в а ле н тнос т и топологич. пространств. В основе этого отношения лежит понятие гомотопии одного непрерывного отображения в другое, имеющее чисто теоретико-множественную природу.



Хотя методы алгебраич. Т. играют исключительно важную роль в топологич. исследовании, существенную роль здесь исполняют и чисто теоретико-множественные конструкции. Это связано с тем, напр., что отношение гомотопич. эквивалентности, примененное к многообразиям, выводит за пределы класса многообразий. При этом получаются более простые топологич. объекты, рассматривать к-рые в технич. отношении весьма полезно. Методы теории гомотопий требуют осуществления теоретико-множественных конструкций типа различных "выметаний», приклеивания одного топологич. пространства к другому вдоль произвольного непрерывного отображения и т. Д. Это приводит к понятиям клеточного разбиения и клеточного пространства; последние и составляют, по-видимому, тот максимально широкий класс пространств, к-рый включает все дифференцируемые многообразия, полиэдры и допускает достаточно полное исследование методами алгебраич. Т.



Для более широких классов пространств - таких, как класс бикомпактов, класс всех паракомпактов или класс всех метризуемых пространств,- ставить проблему различения пространств с точностью до гомеоморфизма посредством обозримой системы вычислимых топологич. инвариантов не представляется возможным ввиду ее интуитивной неразрешимости. На передний план в качестве основной задачи Т. выдвигается здесь задача сравнения не отдельных топологич. пространств, а целых классов топологич. пространств, к-рые, особенно при аксиоматич. подходе, обычно соответствуют отдельным топологич. инвариантам или их комбинациям. При таком подходе основной вопрос Т. превращается в задачу систематич. сравнения топологич. инвариантов. На этом пути удается построить систематическую и развитую классификацию топологич. пространств.



Два метода выступают на первый план в решении этой задачи. Во-первых, метод взаимной классификации пространств и отображений. Речь идет об исследовании поведения топологич. инвариантов при разного типа непрерывных отображениях и о том, когда топологич. пространство из одного данного класса можно представить как образ пространства из другого данного класса при непрерывном отображении того или иного типа. Эта задача тем более важна и естественна, что часто топологич. пространства бывают даны уже будучи связанными нек-рыми непрерывными отображениями напр., когда новое пространство строится как факторпространство нек-рого исходного топологич. пространства.



Второй метод сравнения заключается в применении кардинальных, или кардинальнозначных, топологич. инвариантов, наз. также мацностными характеристиками. Значениями кардинальных инвариантов являются бесконечные кардинальные числа, что дает возможность их сравнивать, пользуясь операциями и законами кардинальных чисел. Это направление Т. выходит на глубокие положения теории множеств - такие, как аксиома Мартина, континуум-гипотеза. На языке кардинальных инвариантов формулируется гипотеза Суслина, неразрешимость к-рой в рамках системы аксиом Цермело - Френкеля теории множеств доказана. Вот характерное рассуждение с кардинальными инвариантами. Для метризуемых пространств плотность и вес совпадают. Значит, если вес и плотность для данного пространства различимы, то оно не метризуемо. В теории кардинальных инвариантов получено много тонких и неожиданных результатов.



Несмотря на отмеченную выше специфику, к-рую приобретают топологич. задачи и методы в зависимости от того, какой класс топологич. пространств выбран для изучения, ряд основных задач, определяющих развитие Т., формулируется общим образом для всех ее разделов и решается на основе нек-рых общих принципов и методов.



К әтим задачам относятся, в частности, следующие: 1) Построение системы топологич. инвариантов на основе $\mathrm{T}$. или внешних структур, ее порождающих. B этом случае возникает задача нахождения этих инвариантов для отдельных пространств и классов пространств.



\begin{enumerate}

  \setcounter{enumi}{1}

  \item Исследование поведения топологич. инвариантов при основных операциях над топологич. пространствами, в частности при переходе к подпространству.



  \item Исследование поведения топологич. инвариантов при разного типа непрерывных отображениях (в частности, вложениях).



  \item Исследование соотношений между топологич. свойствами пространств и их дополнений в нек-ром объемлющем пространстве. Многие результаты геометрич. Т., теоремы двойственности, результаты, связывающие свойства топологич. пространств и их наростов в бикомпактных хаусдорфовых расширениях хорошо иллюстрируют это направление.



\end{enumerate}



K числу общих методов, применяемых в решении большинства задач Т. во всех ее разделах, относятся, в частности, следующие:



\begin{enumerate}

  \item Метод покрытий. Этот метод дает результат при решении проблемы метризации, при определении и исследовании паракомпактных пространств, при определении и исследовании основных объектов комбинаторной Т.- симплициальных и клеточных комплексов. На методе покрытий, в частности на понятии нерва покрытия, основана аппроксимация топологич. пространств полиэдрами. На основе открытых покрытий и отвечающих им разбиений единицы доказываются теоремы о погруякения многообразий в евклидово пространство.



  \item Метод функторов. Он заключается в сопоставлении топологич. пространствам алгебраич. и алгебротопологич. объектов, обладающих правильным - функториальным - поведением и допускающих вычисление. Гомологий группа, когомологий кольцо, $K$-функтор, связанный с понятием векторного расслоения над топологич. пространством, обобщающим понятие касательного многообразия, - важные примеры функторов. На применении таких функторов основан алгебраич метод в $\mathrm{T}$.



  \item Метод спектров. Суть его заключается в представлении более сложно устроенных пространств в качестве предела обратного спектра из пространств более простых (напр., полиэдров), при этом изучается связь между топологич. инвариантами элементов спектра и



\end{enumerate}



\end{document}